\chapter{Couche 4 — Le pilotage du parc}

Les machines existent, mais rien ne décide quand les créer, où les placer, ni
quand les détruire. Sept briques pour transformer un démon Docker en un service
de location.

% -----------------------------------------------------------------------------
\outil{subprocess}{brique 22 · piloter Docker depuis Python}

\quoi
Le module de la bibliothèque standard qui lance un programme externe et récupère
sa sortie. Ici, il lance la commande \cmd{docker}.

\pourquoi
Il existe une bibliothèque Python officielle, \cmd{docker-py}, qui parle à
l'API Docker sans passer par la ligne de commande. Elle est plus élégante.

On l'a écartée pour une raison de \textbf{conformité et de lisibilité}. Le sujet
impose une commande précise :

\begin{center}
\cmd{docker run -d --restart=always -p <port>:22 <image>}
\end{center}

Avec un SDK, cette commande n'apparaît nulle part : elle devient un appel de
méthode avec des paramètres nommés. Avec \cmd{subprocess}, chaque option se lit
telle quelle dans le code, et l'on peut montrer à la prof la ligne exacte qui
part.

\role
Une seule fonction centralise tous les appels, dans
\fichier{2\_webapp/docker\_ops.py} :

\begin{lstlisting}[language=Python]
def run_docker(*args, worker="local", timeout=DOCKER_TIMEOUT) -> str:
    cmd = [*wk.docker_command(worker), *args]
    # Appel sans shell : les arguments sont passes en liste, donc jamais
    # re-interpretes par un interpreteur de commandes.
    res = subprocess.run(cmd, capture_output=True, text=True, timeout=timeout)
\end{lstlisting}

\begin{piege}{\texttt{shell=True} est une porte ouverte}
Écrire \cmd{subprocess.run(f"docker run ... \{nom\}", shell=True)} laisserait un
utilisateur nommer sa machine \cmd{x; rm -rf /} et faire exécuter n'importe quoi
au serveur. C'est une \emph{injection de commande}.

Passer une \textbf{liste} d'arguments, sans shell, supprime le problème à la
racine : il n'y a aucun interpréteur pour relire la chaîne. Chaque élément de la
liste est un argument, point.
\end{piege}

\apport
La commande imposée est visible et vérifiable. Et les tests peuvent espionner
cette unique fonction pour vérifier ce qui part, sans Docker.

% -----------------------------------------------------------------------------
\outil{Les étiquettes Docker}{brique 23 · l'identité des machines}

\quoi
Des paires clé-valeur attachées à un conteneur à sa création, et sur lesquelles
on peut filtrer.

\pourquoi
C'est la brique née d'un incident. Le Faucheur doit savoir quels conteneurs lui
appartiennent. La première version les reconnaissait à leur nom, qui commence
par \cmd{insacloud\_}.

Puis la pile \cmd{docker compose} est arrivée, avec des conteneurs nommés
\cmd{insacloud-web}, \cmd{insacloud-db}, \cmd{insacloud-faucheur}. Le Faucheur
les a pris pour des machines louées orphelines et les a détruits.
\textbf{Il s'est supprimé lui-même.}

\role
Deux étiquettes, apposées à toute machine louée :

\begin{lstlisting}[language=Python]
MACHINE_LABEL = "insacloud.role"          # = "rented-machine"
CLUSTER_LABEL = "insacloud.cluster"       # = nom du deploiement
\end{lstlisting}

Les deux filtres sont combinés : une machine doit être une machine louée
\textbf{et} appartenir à ce déploiement. Le second a été ajouté après un
deuxième incident : deux déploiements partageant un démon se détruisaient
mutuellement leurs conteneurs.

\begin{retenir}
Un nom est une commodité d'affichage ; une étiquette est une identité. Dès
qu'un programme doit décider de \textbf{détruire} quelque chose, il doit se
fonder sur une marque qu'il a lui-même apposée, jamais sur une convention de
nommage.
\end{retenir}

\apport
Le Faucheur ne peut plus se tromper de cible. Deux plateformes peuvent
cohabiter sur un même démon Docker.

% -----------------------------------------------------------------------------
\outil{Le transport SSH de Docker}{brique 24 · piloter les nœuds distants}

\quoi
Le client Docker sait parler à un démon distant à travers une connexion SSH :
\cmd{docker -H ssh://utilisateur@machine ps}.

\pourquoi
L'autre façon de commander un démon distant est d'exposer son API sur le réseau,
port 2375 en clair ou 2376 en TLS. C'est plus direct, et c'est une très mauvaise
idée : \textbf{qui parle au démon Docker d'une machine en est root}. Une erreur
de pare-feu, et toute la plateforme est compromise.

Le transport SSH ne demande aucun port supplémentaire, réutilise une
authentification par clé déjà éprouvée, et chiffre tout.

\role
Une fonction construit le préfixe de commande selon le nœud visé :

\begin{lstlisting}[language=Python]
def docker_command(worker: str = "local") -> list:
    host = worker_host(worker)
    if host is None:
        return ["docker"]
    return ["docker", "-H", f"ssh://{DOCKER_SSH_USER}@{host}"]
\end{lstlisting}

Conséquence heureuse : \textbf{le code est le même} en mono-hôte et en cluster.
Seul le préfixe change. C'est ce qui permet de développer sur un portable sans
monter quatre machines virtuelles.

\apport
Un parc distribué sans ouvrir un seul port Docker sur le réseau.

% -----------------------------------------------------------------------------
\outil{L'agent des nœuds}{brique 25 · l'inscription spontanée}

\quoi
Un programme Python de 263 lignes qui tourne sur chaque worker. Au démarrage il
se présente au contrôleur ; ensuite il signale sa présence toutes les trente
secondes.

\pourquoi
Sans lui, il faudrait écrire la liste des workers dans la configuration du
contrôleur. Ajouter un nœud deviendrait : modifier un fichier, redéployer. Et un
nœud en panne resterait dans la liste jusqu'à ce qu'un humain l'en retire.

L'agent inverse la responsabilité : c'est le nœud qui se déclare. Le contrôleur
ne contient \textbf{aucune adresse de worker}.

Il n'utilise que la bibliothèque standard --- \cmd{urllib} pour les appels,
\cmd{/proc/meminfo} pour la mémoire --- précisément pour qu'il n'y ait rien à
installer sur les workers.

\role
Deux routes de l'API : \cmd{POST /workers/register} et
\cmd{POST /workers/heartbeat}, protégées par un jeton partagé comparé en temps
constant. L'inscription est retentée avec un recul exponentiel : au démarrage
d'un cluster, un worker peut être prêt avant le contrôleur.

\begin{center}
\small
\begin{tabularx}{\textwidth}{@{}lX@{}}
\toprule
\textbf{État} & \textbf{Signification} \\
\midrule
\cmd{AVAILABLE} & le nœud bat, Docker répond, il reste de la place \\
\cmd{BUSY} & le nœud bat, mais il est plein : on ne lui envoie plus rien \\
\cmd{OFFLINE} & plus de battement depuis 90 s : ses machines sont recréées ailleurs \\
\bottomrule
\end{tabularx}
\end{center}

\begin{retenir}
La distinction \cmd{BUSY} / \cmd{OFFLINE} est essentielle. Un nœud saturé est en
bonne santé : migrer ses machines serait coûteux et inutile. Confondre « plein »
et « mort » déclencherait des reprises à répétition.
\end{retenir}

\apport
Pour ajouter un worker, on l'allume. C'est tout. C'est la Fonctionnalité 7 du
sujet, et c'est ce qui rend le parc élastique.

% -----------------------------------------------------------------------------
\outil{L'ordonnanceur}{brique 26 · choisir le nœud}

\quoi
La fonction qui décide sur quel worker placer une nouvelle machine.

\pourquoi
Les stratégies possibles sont nombreuses : tourniquet, aléatoire, par
pondération de ressources. On a choisi la plus simple qui soit défendable ---
le nœud joignable le \textbf{moins chargé} --- pour une raison précise : elle
est \textbf{déterministe}, donc testable.

\begin{retenir}
Un ordonnanceur dont le départage en cas d'égalité est aléatoire donne un test
intermittent. Et un test intermittent finit toujours par être ignoré. Ici,
l'égalité se tranche par ordre alphabétique.
\end{retenir}

\role
Trois règles : ne considérer que les nœuds \cmd{AVAILABLE} ; parmi eux, prendre
celui qui héberge le moins de machines ; départager par le nom. Puis réserver un
port libre, dont l'unicité est garantie par l'index partiel de la couche 2.

\apport
La charge se répartit, et le comportement est reproductible --- donc
vérifiable par un test.

% -----------------------------------------------------------------------------
\outil{Le Faucheur}{brique 27 · le démon de nettoyage}

\quoi
Un processus séparé qui tourne en boucle, toutes les dix secondes, avec trois
missions : détruire ce qui a expiré, réconcilier Docker et la base, et
surveiller le parc.

\pourquoi
Personne ne viendra cliquer sur « supprimer » à trois heures du matin. Il faut
un service qui surveille les horloges.

Pourquoi un \textbf{processus séparé} plutôt qu'un fil d'exécution dans
l'application ? Parce que Gunicorn lance quatre processus : il y aurait quatre
Faucheurs qui se marcheraient dessus. Un service distinct, démarré une fois par
systemd, évite le problème et peut être redémarré sans couper le site.

\role
Trois tâches à chaque cycle.

\textbf{Expirer} : une seule requête SQL, celle que sert l'index
\cmd{(status, end\_time)}.

\textbf{Réconcilier} : la base dit ce qui \emph{devrait} exister, Docker ce qui
existe \emph{vraiment}. Un conteneur orphelin est supprimé ; une instance sans
conteneur passe à « arrêtée ».

\textbf{Surveiller le parc} : déclarer hors ligne les nœuds muets, et faire
recréer leurs machines ailleurs.

\begin{piege}{Un nœud injoignable n'est pas un nœud vide}
C'est le piège le plus dangereux de la réconciliation. Si le nœud ne répond pas,
\cmd{docker ps} renvoie une liste vide. Un code naïf en conclurait que toutes
ses machines ont disparu et les marquerait « arrêtées » --- alors qu'elles
tournent parfaitement.

D'où la règle : un nœud injoignable est \textbf{ignoré}, jamais interprété.
\end{piege}

\apport
Les machines disparaissent toutes seules à l'échéance, les ports se libèrent, et
la base ne dérive plus de la réalité. C'est la Fonctionnalité 12, et la moitié
de la 13.

% -----------------------------------------------------------------------------
\outil{systemd}{brique 28 · faire tourner les services, et les réparer}

\quoi
Le gestionnaire de services de Linux : il démarre les programmes au bon moment,
les redémarre s'ils meurent, et collecte leurs journaux.

\pourquoi
L'alternative serait de lancer les processus à la main dans un \cmd{screen}, ou
d'écrire des scripts d'init. systemd apporte trois choses qu'on ne veut pas
réécrire : le redémarrage automatique, l'ordre de démarrage selon les
dépendances, et les journaux consultables par \cmd{journalctl}.

\role
Deux unités et un \emph{timer}.

L'unité de l'application lance Gunicorn sous un compte sans privilège, avec les
secrets lus dans un fichier séparé :

\begin{lstlisting}
[Service]
User=insacloud
EnvironmentFile=/etc/insacloud/insacloud.env     # secrets, en 0600
ExecStart=/opt/insacloud/.venv/bin/gunicorn --preload -w 4 ...
Restart=always
\end{lstlisting}

\begin{piege}{Un secret dans une unité systemd est public}
Mettre \cmd{Environment=INSACLOUD\_VAULT\_KEY=...} dans le fichier d'unité le
rend lisible par tous : les unités sont en lecture pour tout le monde, et
\cmd{systemctl show} les affiche. La bonne façon est un fichier séparé en
\cmd{0600}, référencé par \cmd{EnvironmentFile=}.
\end{piege}

\textbf{Le watchdog} est un \emph{timer} qui se déclenche chaque minute. Il
vérifie que nginx répond, que le site renvoie bien du HTTPS, et que le Faucheur
tourne. Il redémarre ce qui est tombé et journalise les workers injoignables.

Pourquoi un timer systemd plutôt que cron ? Parce qu'il s'intègre aux journaux,
connaît l'état du service qu'il surveille, et ne démarre pas avant que ses
dépendances soient prêtes. Cron ne sait rien de tout cela.

\apport
Le service se relance seul après une panne ou un redémarrage du serveur. C'est
le troisième ajout au-delà du sujet.

% -----------------------------------------------------------------------------
\section*{Ce que la plateforme sait faire après la couche 4}

Elle loue, place, surveille et détruit. Elle survit à la panne d'un nœud. Elle
se répare elle-même. Mais tout cela n'existe que sur des machines que
quelqu'un a installées à la main --- et ça, ce n'est pas tenable.
